%% ============================================================
%%  SOLUCIÓN — Ejercicio 2.3: Sistema de ecuaciones con align
%%  Clase 1 — Después de diapositiva "Ecuaciones alineadas: align"
%%  Compilar: F5 (pdflatex, dos veces)
%% ============================================================

\documentclass[12pt, a4paper]{article}
\usepackage[utf8]{inputenc}
\usepackage[T1]{fontenc}
\usepackage[spanish]{babel}
\usepackage{amsmath, amssymb}

\title{Derivación del estimador OLS}
\author{Tu Nombre}
\date{\today}

\begin{document}
\maketitle

\section{Demostración paso a paso}

El estimador OLS se obtiene minimizando la suma de cuadrados de residuos:

\begin{align}
  \hat{\beta} &= \arg\min_{\beta} \| Y - X\beta \|^2 \\
              &= \arg\min_{\beta} (Y - X\beta)'(Y - X\beta) \nonumber \\
              &= (X'X)^{-1}X'Y \label{eq:ols_final}
\end{align}

La interpretación de estos pasos es la siguiente:

\begin{enumerate}
  \item En la línea \eqref{eq:ols_final}, expresamos la suma de cuadrados usando 
        la norma euclidiana.
  \item Expandimos el producto interno.
  \item Tomamos la derivada con respecto a $\beta$, igualamos a cero, y resolvemos 
        (pasos algebraicos omitidos).
  \item Obtenemos la fórmula cerrada para el estimador OLS.
\end{enumerate}

\subsection*{Interpretación de $\|Y - X\beta\|^2$}

Este símbolo representa la norma euclidiana (distancia en $\mathbb{R}^N$) del 
vector de residuos. La norma es la raíz cuadrada de la suma de cuadrados, así 
que al elevar al cuadrado obtenemos precisamente $\sum_{i=1}^N (Y_i - X_i'\beta)^2$.

\end{document}
